\section{问题重述}
\subsection{问题背景}
小区用电需要连续供给，而光伏出力随时间变化，外网购电价格也有峰谷之分。储能可以把富余光伏或低价购电转移到其他时段使用，但能量损耗、容量和功率限制使这种转移不能无限进行。因此，微网调度的任务是在满足负载的前提下，联合安排外网购电和电池充放电，使实际购电费用尽可能低。

题目提供分时电价、负载与光伏数据、定时发布的光伏预报及储能设备参数\cite{problem}。储能额定容量为12000 kWh，允许运行范围为1200至10800 kWh，最大充放电功率为5000 kW，初始储电量为6000 kWh。围绕这些条件，本文研究以下三个问题。

\subsection{需要解决的问题}
\begin{enumerate}[label=(\arabic*),leftmargin=2.2em]
\item 在已知一天电价、负载和光伏预测的条件下，于0:00确定全天购电与储能调度，满足日初、日末储电量相同，给出指定时段购电量、分段充放电量和全天费用。
\item 在负载与光伏逐日变化、当天购电计划必须提前锁定的条件下，利用历史数据制定计划，并通过储能与5倍电价的紧急购电补足实时缺口，给出2月至12月的连续运行结果。
\item 在0:00、6:00、12:00和18:00获得更新光伏预报的条件下，考虑取消及新增购电的结算费用，修订当天剩余计划，并判断采用各次日内预报是否具有经济价值。
\end{enumerate}

\section{问题分析}
\subsection{共同的决策基础}
三个问题的物理对象相同：购电和光伏首先满足负载，富余电量可进入储能；供给不足时由储能放电或外网补电。储电量把不同时段连接起来，因此不能仅按某一时刻的电价作局部选择。例如，低价时段充得过满，可能挤占随后光伏的存储空间；过早放电，又可能增加晚高峰的购电需求。

将一天离散为144个十分钟区间后，供需平衡、储电递推和设备边界均可写为线性关系。本文据此建立统一的线性规划框架，再根据各问的信息条件与结算规则扩展决策。

\subsection{各问的递进关系}
问题一的主要矛盾是能量在不同时段之间如何转移。给定全天输入，可以一次性优化购电与储能轨迹，并通过无储能对照解释节省费用的来源。

问题二增加了预测误差。提前多购的电量即使没有利用也要付费，而少购可能引起高价紧急补电。本文用历史误差构造若干可能的负载与光伏情景，让各情景共用一份日前购电计划，再根据实际观测滚动安排储能。这延续了两阶段随机规划中“先决策、后补救”的建模思想\cite{stochastic}。

问题三增加了利用新信息修订购电的机会。是否调整取决于预计缺口、当前储电量和交易成本，而不只取决于预报是否发生变化。采用滚动优化思想，每次用最新可用信息计算未来安排，只执行当前动作，随后根据新状态再次优化\cite{mpc}。其中购电修订仅发生在允许的发布时间，储能则每十分钟更新。

\begin{figure}[htbp]\centering
\includegraphics[width=.98\textwidth,height=.40\textheight,keepaspectratio]{figures/overview_flow.pdf}
\caption{日前计划、日内购电修订与实时储能执行的关系。日内购电修订仅适用于问题三。}\label{fig:flow}
\end{figure}

\section{模型假设}
\begin{enumerate}[label=(\arabic*),leftmargin=2.2em]
\item 每个十分钟区间内负载、实测光伏功率及电价按常值处理，当前区间功率可即时测得；控制器只能使用当前及此前的实测数据和已经发布的预报。
\item 充电和放电单程效率均为0.9，最大充放电功率按微网侧定义。不计未给定的自放电、寿命折损和启停费用。
\item 储能容量与功率严格满足题目边界；外网能够补足供电缺口，不另设题目未给定的外网购电功率上限。不考虑售电收入，富余供给允许舍弃。
\item 问题一采用6000 kWh的日初储电量，并在日末恢复。问题二、三以1月为历史学习期，电池保持6000 kWh，从2月1日起连续执行至12月31日，实际储电量跨日传递。
\item 问题二、三在每日0:00采用48小时规划视野，日内规划终点保持为次日24:00，末端规划储电量为6000 kWh。次日购电只作前瞻，次日0:00重新确定；实际每日午夜不强制恢复6000 kWh。
\item 问题二的计划电量全额收费，未利用部分不退款；紧急购电按当时电价的5倍结算。问题三取消电量退回原购电费并收取50\%违约费，新增电量按1.5倍电价结算，均相对0点原计划按最终交付量计算一次。
\end{enumerate}

\section{符号说明}
\begin{table}[htbp]\centering\small
\caption{主要符号及单位}\label{tab:symbols}
\begin{tabular}{clc}\toprule
符号 & 含义 & 单位\\\midrule
$t,d,s$ & 时段、日期和情景索引 & --\\
$\Delta t,T$ & 时段长度$1/6$小时、每日区间数144 & h，--\\
$p_t$ & 时段普通购电价格 & 元/kWh\\
$\ell_t,v_t$ & 负载电量、光伏电量 & kWh\\
$\widehat\ell_t,\widehat v_t$ & 负载和光伏预测电量 & kWh\\
$g_t,a_t$ & 0点原计划购电量、最终有效购电量 & kWh\\
$q_t$ & 统一模型中的有效普通购电供给 & kWh\\
$c_t,d_t$ & 微网侧充电输入、放电输出电量 & kWh\\
$e_t,w_t$ & 紧急购电量、未利用供给电量 & kWh\\
$E_t$ & 第$t$区间末储电量，$E_0$为初值 & kWh\\
$\eta,\pi_s$ & 单程效率、情景概率 & --\\
$\delta_t^+,\delta_t^-$ & 相对原计划的最终增购、减购电量 & kWh\\
$C,J$ & 实际费用、规划目标值 & 元\\
\bottomrule\end{tabular}
\end{table}
日期索引在需要区分跨日状态时写作$E_{d,t}$；放电变量$d_t$的下标始终表示时段。情景变量采用上标$(s)$。除专门标明功率的输入和图形外，模型中的购电、充放电及负载均使用区间电量，避免重复乘以区间长度。

\section{数据处理与特征分析}
\subsection{时间统一与数据用途}
将一天划分为$[0{:}00,0{:}10),\ldots,[23{:}50,24{:}00)$，每段功率乘以$\Delta t=1/6$得到kWh电量。附件一提供问题一输入，并为后两问提供每天重复的分时电价；附件二的实测负载与光伏用于重建历史预测及逐段回测；附件三的发布预报仅在问题三中使用。

后两问在1月积累历史信息，以2月1日至12月31日共334天作为统一评价区间。预测误差按历史上真实可见的信息重建，不用全年数据先拟合误差后再回测。光伏为零时仍保留原始零值，误差分析采用电量差，避免夜间零出力导致相对误差无定义。

附件三每次发布的24个功率点分别对应发布后第1至第24小时。以发布时刻刚结束区间的已知功率作为起点，相邻整点之间线性插值，并在十分钟区间上积分得到预测电量。新预报覆盖不到的远端时段，以最近完整一天的对应光伏数据补充前瞻。

\subsection{输入特征与调度含义}
定义净负荷$n_t=\ell_t-v_t$。净负荷为正时，需要外网或电池提供差额；净负荷为负时，光伏具有剩余，是否可以保留这部分能量取决于电池的剩余容量与充电功率。图\ref{fig:inputs}将净负荷与电价分面展示，便于同时观察能量缺口和购电时机。

\begin{figure}[htbp]\centering
\includegraphics[width=.96\textwidth,height=.47\textheight,keepaspectratio]{figures/overview_inputs.pdf}
\caption{附件一的电价、负载、光伏与净负荷。功率和电价分别使用独立纵轴。}\label{fig:inputs}
\end{figure}

光伏出力集中在白天，负载与光伏并不逐时匹配，因此即使全天发电总量较高，局部时段仍可能缺电。储能能够缓解这种错配，但往返损耗也会增加为负载交付同等电量所需的输入。故本文既比较费用，也核对未利用供给、储能损耗和期末储电量，避免仅凭购电总量减少判断方案优劣。

\section{模型建立与求解}
\subsection{统一的能量平衡与储能约束}\label{sec:common}
记$q_t$为当前有效的普通购电供给量：问题一、二取$q_t=g_t$，问题三取$q_t=a_t$。共同物理模型为
\begin{align}
q_t+v_t+d_t+e_t&=\ell_t+c_t+w_t,\label{eq:balance}\\
E_t&=E_{t-1}+\eta c_t-d_t/\eta,\label{eq:soc}\\
1200\le E_t&\le10800,\qquad
0\le c_t,d_t\le5000/6,\label{eq:limits}\\
q_t,e_t&\ge0,\qquad 0\le w_t\le q_t+v_t.\label{eq:nonnegative}
\end{align}
式\eqref{eq:balance}保证负载获得足额供电；式\eqref{eq:soc}中只有充放电改变电池电量，未利用供给不会扣减储电量。问题一进一步取$e_t=0$、$w_t\le v_t$。问题二、三的$w_t$还可包括已经支付但没有利用的计划电量。

微网输入100 kWh只增加90 kWh储电量，该部分再释放时能向微网提供81 kWh。这个损耗关系决定了低价充电并不总是经济。实际执行只允许充电或放电中的一种；具体的当前供需规则与线性规划等价解处理见附录\ref{app:implementation}。

各模型的目标与上述约束均为线性或可线性化的分段线性形式，使用HiGHS求解。相关单纯形算法见文献\cite{simplex}。最优性针对给定预测和情景的规划问题，在线运行效果则由真实序列上的实际账单评价。
